\renewcommand{\abstractname}{Abstract} % Veränderter Name für das Abstract
\begin{abstract}
\begin{addmargin}[1.5cm]{1.5cm}        % Erhöhte Ränder, für Abstract Look
\thispagestyle{plain}                  % Seitenzahl auf der Abstract Seite

\begin{center}
\small\textit{- Deutsch -}             % Angabe der Sprache für das Abstract
\end{center}

\vspace{0.25cm}

Die Vorhersage von Aktienkursen stellt aufgrund der Volatilität und nicht-linearen Natur von Aktienkursen eine große Herausforderung dar. Diese Arbeit stellt eine vollständige Pipeline zur algorithmischen Aktienkursvorhersage und Evaluation vor. Aus dem S\&P 500 werden mittels stochastischer Differentialgleichungen unter Berücksichtigung makroökonomischer Faktoren zehn profitable Aktien ausgewählt. Daraufhin werden drei unterschiedliche Modellansätze zur Aktienkursvorhersage untersucht. LinearSeparation ist ein neu entwickeltes Verfahren basierend auf HP-Filtern und ARMA, Preprocessed SVR verwendet Support Vector Regression mit technischen Indikatoren und Sentiment-Daten und TimesFM verwendet ein vortrainiertes Foundation Model im Zero-Shot-Modus. Die Modelle werden auf den zehn gewählten Aktien getestet und evaluiert.
Die Evaluation erfolgt zweistufig. Zunächst werden die Modelle anhand statistischer Fehlermaße (MAPE, RMSE, OOS-$R^2$, Directional Accuracy) bewertet und anschließend in einer simulierten Handelsstrategie mit Half-Kelly-Positionsgrößen und signalbasiertem Trading getestet. Die Ergebnisse zeigen, dass trotz guter Punktvorhersagequalität (MAPE 1,08--1,30\%) keines der Modelle die Buy-and-Hold-Strategie konsistent übertrifft. SVR erzielt mit durchschnittlich 6,56\% Return und dem niedrigsten Maximum Drawdown (13,95\%) die beste Performance. Die Diskrepanz zwischen präzisen Punktvorhersagen und schwacher Trading-Performance verdeutlicht die fundamentale Schwierigkeit, aus guten Vorhersagen profitable Handelssignale abzuleiten.


\end{addmargin}
\end{abstract}